\chapter{A Concrete Naming Certificate for $\SeparatedMetricCompletionUp$}
\label{ch:concrete-instances-separated-metric-completion-namecert}
\origin{human}

Following Bishop--Bridges and Kelley, the $\SeparatedMetricCompletionUp$ packet names the finite completion boundary obtained after zero-distance identifications are made visible rather than quotient-hidden. It sits between the metric completion source of \autoref{ch:concrete-instances-metriccompletion-namecert}, the separated metric classifier of \autoref{ch:concrete-instances-separatedmetric-namecert}, the Cauchy-filter route of \autoref{ch:concrete-instances-cauchyfilter-namecert}, the Cauchy-net route of \autoref{ch:concrete-instances-cauchynet-namecert}, the regular rational readback of \autoref{ch:concrete-instances-regseqrat-namecert}, and the terminal real-name seal of \autoref{ch:concrete-instances-real-namecert}. The packet exposes a separated completion handoff for L10 Real and Completion consumers; it is not a quotient completion, selected representative principle, host equality theorem, countable-choice selector, or ambient completeness assertion.

\begin{definition}[SeparatedMetricCompletion carrier]
\label{def:separated-metric-completion-carrier}
A $\SeparatedMetricCompletionUp$ carrier is a finite $\BHist$ packet
$$
S=(M,B,Q,R,Z,H,C,P,N)
$$
with the following displayed rows:
\begin{enumerate}
\item $M$ is the $\MetricCompletionUp$ source row naming the completion request and its metric evidence.
\item $B$ is the selected Cauchy branch row, either the $\CauchyFilterUp$ route or the $\CauchyNetUp$ route, together with its displayed observation schedule.
\item $Q$ is the $\RegSeqRatUp$ and $\RealUp$ readback handoff extracted from the same finite Cauchy branch.
\item $R$ records the $\SeparatedMetricUp$ zero-distance classifier used to identify when two completion reads name the same separated point.
\item $Z$ is the local zero-distance comparison ledger tying $R$ to the source metric row $M$ and the branch row $B$.
\item $H$ records componentwise $\hsame$ transport for the metric-completion source, Cauchy branch, readback, zero-distance classifier, comparison ledger, continuation, provenance, and local naming rows.
\item $C$ records displayed $\Cont$ replay from a completion request through the selected branch, regular readback, Real seal, zero-distance classifier, and separated handoff.
\item $P$ is the $\Pkg$ provenance over $\MetricCompletionUp$, $\SeparatedMetricUp$, $\CauchyFilterUp$, $\CauchyNetUp$, $\RegSeqRatUp$, $\RealUp$, $\BHist$, $\hsame$, $\Cont$, and $\NameCert$.
\item $N$ is the local $\NameCert_{\SeparatedMetricCompletionUp}$ row.
\end{enumerate}
The carrier contains no coordinate for quotient completion classes, selected representatives, host equality, countable choice, ambient completeness, or an unrecorded completion object.
\end{definition}

\begin{theorem}[SeparatedMetricCompletion NameCert obligations]
\label{thm:separated-metric-completion-namecert-obligations}
For every accepted $\SeparatedMetricCompletionUp$ carrier $S=(M,B,Q,R,Z,H,C,P,N)$, the local $\NameCert_{\SeparatedMetricCompletionUp}$ surface consists exactly of these obligation rows:
\begin{enumerate}
\item carrier admission for the displayed $\MetricCompletionUp$ source, selected Cauchy branch, regular readback, Real seal, zero-distance classifier, comparison ledger, transport, continuation, provenance, and local naming rows;
\item source-first exposure, requiring every separated completion read to consume $M$ before it reads the branch row $B$ or the readback row $Q$;
\item branch exactness, requiring $B$ to be one displayed $\CauchyFilterUp$ or $\CauchyNetUp$ route and requiring $Q$ to be read from that same route;
\item separation exactness, requiring the zero-distance classifier $R$ and comparison ledger $Z$ to mediate all same-point reads before a consumer treats the handoff as separated;
\item ledger refusal, excluding quotient completion classes, selected representatives, host equality, countable-choice selectors, ambient completeness, and detached completion objects.
\end{enumerate}
No local NameCert obligation is stored outside those rows.
\end{theorem}
\begin{proof}
Project the carrier of \autoref{def:separated-metric-completion-carrier}. Its public rows are exactly $M,B,Q,R,Z,H,C,P,N$, so carrier admission exposes a metric-completion source, one Cauchy branch, the readback handoff, the zero-distance classifier, and the structural rows.

The source and branch obligations follow by replaying $C$. A separated completion read first consumes $M$, then chooses the displayed branch row $B$, and only then reads the $\RegSeqRatUp$ and $\RealUp$ handoff in $Q$.

The separation obligation is carried by $R$ and $Z$. These rows are the only coordinates that compare completion reads as zero-distance data, while $H$, $P$, and $N$ transport, package, and name the same displayed route.

Every excluded object is absent from the carrier. Quotient classes, selected representatives, host equality, countable choice, ambient completeness, and detached completion objects therefore cannot supply an additional local obligation.
\end{proof}

\begin{theorem}[SeparatedMetricCompletion MetricCompletion handoff]
\label{thm:separated-metric-completion-metriccompletion-handoff}
Let $S=(M,B,Q,R,Z,H,C,P,N)$ be an accepted $\SeparatedMetricCompletionUp$ packet. Every $\MetricCompletionUp$ consumer that reads the separated handoff factors through the finite route
$$
M \longmapsto B \longmapsto Q \longmapsto R,Z \longmapsto H,C,P,N .
$$
Thus the completion source is read before the Cauchy-filter or Cauchy-net branch, the Real-facing readback is read from that same branch, and the $\SeparatedMetricUp$ zero-distance classifier mediates separated equality before the packet exits to downstream L10 consumers.
\end{theorem}
\begin{proof}
Start with an accepted separated completion read. By carrier admission, the only completion source available to the consumer is the displayed row $M$.

Replay the selected branch. The row $B$ records whether the finite Cauchy evidence is filter-presented or net-presented, and the readback row $Q$ is attached to that same branch through the continuation row $C$.

Now apply the separated classifier rows. The row $R$ supplies the $\SeparatedMetricUp$ zero-distance classifier, and $Z$ records its compatibility with the source metric and the selected branch. Hence same-point use is mediated by displayed zero-distance data rather than by host equality.

Finally pass through the structural rows $H,C,P,N$. They preserve and name the route already displayed, so the handoff cannot introduce quotient completion classes, selected representatives, countable choice, ambient completeness, or a detached completion object.
\end{proof}

\begin{closurestatus}{\SeparatedMetricCompletionUp}
  \theoryclosure{\seedClosure}
  \scopeclosed{\autoref{def:separated-metric-completion-carrier}, \autoref{thm:separated-metric-completion-namecert-obligations}, and \autoref{thm:separated-metric-completion-metriccompletion-handoff} seed the finite separated completion boundary between metric completion evidence, Cauchy-filter or Cauchy-net presentation, regular readback, Real sealing, and zero-distance classification.}
  \formalstatus{\unformalizedV}
  \bridgestatus{none}
  \notclaimed{This chapter does not close quotient completion classes, selected representatives, host equality, countable-choice selectors, ambient completeness, Lean verification, or bridge checking.}
  \upgradepath{Obligation closure requires carrier admission, branch exactness, separation exactness, MetricCompletion handoff, and ledger refusal proofs for the same packet.}
  \constructivestory{A $\SeparatedMetricCompletionUp$ packet is generated from a MetricCompletionUp source row, a displayed CauchyFilterUp or CauchyNetUp branch, RegSeqRatUp readback, a RealUp seal, a SeparatedMetricUp zero-distance classifier, componentwise hsame transport, displayed Cont replay, Pkg provenance, and a local NameCert row.}
\end{closurestatus}
